\documentclass[a4paper,twocolumn]{jsarticle}
% a4jでは画像挿入時に問題が起こる場合があるため，a4paperを使用する．

% ===== 基本 =====
\usepackage[
    a4paper,
    left=2.0cm,
    right=2.0cm,
    top=2.0cm,
    bottom=2.0cm
]{geometry}
\usepackage{amsmath,amssymb,mathtools,bm}
\usepackage[dvipdfmx]{graphicx}
\usepackage[dvipdfmx]{xcolor}
\usepackage[dvipdfmx,unicode]{hyperref}
\usepackage{pxjahyper}
\usepackage{titlesec}
\usepackage{wrapfig}
\usepackage{multicol}
\usepackage{enumitem}
\usepackage{listings}
\usepackage{float}
\usepackage{caption}
\usepackage{subcaption}

% ===== TikZ・図 =====
\usepackage{tikz}
\usetikzlibrary{positioning}
\usepackage{pgfplots}
\pgfplotsset{compat=1.18}

% ===== 化学 =====
\usepackage[version=4]{mhchem}
\usepackage{chemfig}

% ===== 回路 =====
\usepackage{circuitikz}

% ===== アルゴリズム =====
\usepackage{algorithm}
\usepackage{algpseudocode}

% ===== ダミーテキスト =====
% 発展例で図の回り込みを確認するために使用する．実際の文書では不要．
\usepackage{lipsum}

% ===== 数学補助 =====
% \numberwithin{equation}{section}
\renewcommand{\theequation}{\arabic{equation}}

\DeclareMathOperator{\sgn}{sgn}
\DeclareMathOperator{\tr}{tr}
\DeclareMathOperator{\diag}{diag}

\newcommand{\dif}{\mathrm{d}}
\newcommand{\ee}{\mathrm{e}}
\newcommand{\ii}{\mathrm{i}}

\newcommand{\pd}[2]{\frac{\partial #1}{\partial #2}}
\newcommand{\pdd}[2]{\frac{\partial^2 #1}{\partial #2^2}}
\newcommand{\dd}[2]{\frac{\mathrm{d} #1}{\mathrm{d} #2}}

\newcommand{\vect}[1]{\bm{#1}}
\newcommand{\tensor}[1]{\bm{#1}}

% 画像がまだ配置されていない場合も，文書全体はコンパイルできるようにする．
% 指定した画像が存在すれば，通常どおりその画像を表示する．
\newcommand{\practiceimage}[2][\linewidth]{%
  \IfFileExists{#2}{%
    \includegraphics[width=#1]{#2}%
  }{%
    \fbox{%
      \parbox[c][28mm][c]{\dimexpr#1-2\fboxsep-2\fboxrule\relax}{%
        \centering\small
        画像ファイル\\未配置%
      }%
    }%
  }%
}

% ===== 見出し設定 =====
\titleformat{\section}[block]
  {\normalfont\Large\bfseries}
  {\thesection}{1em}{}

\titleformat{\subsection}[block]
  {\normalfont\large\bfseries}
  {\thesubsection}{1em}{}

\titleformat{\subsubsection}[block]
  {\normalfont\normalsize\bfseries}
  {\thesubsubsection}{1em}{}

\titlespacing*{\section}{0pt}{2.0ex plus 0.5ex minus 0.2ex}{0.8ex}
\titlespacing*{\subsection}{0pt}{1.5ex plus 0.4ex minus 0.2ex}{0.5ex}

\lstset{
    language=Python,
    basicstyle=\ttfamily\small,
    numbers=left,
    numberstyle=\tiny\color{gray},
    frame=single,
    rulecolor=\color{black},
    backgroundcolor=\color{gray!10},
    breaklines=true,
    showstringspaces=false,
    columns=fullflexible,
    keywordstyle=\color{blue},
    commentstyle=\color{green!50!black},
    stringstyle=\color{red},
    tabsize=4
}

\setlist[itemize]{topsep=4pt,itemsep=2pt,leftmargin=*}
\setlist[enumerate]{topsep=4pt,itemsep=2pt,leftmargin=*}
\captionsetup{font=small}
\setlength{\columnsep}{8mm}
\setlength{\emergencystretch}{1em}
\raggedbottom

\title{自然言語処理を用いた\\映画レビューのネタバレ検出モデルの構築}
\author{伊藤 大惺}
\date{\today}

\begin{document}

\maketitle

\begin{abstract}
本稿では，映画レビューにおけるネタバレ（Spoiler）を自動的に検出するモデルの構築について述べる．自然言語処理技術を用いることで，文章の特徴を抽出し，未鑑賞のユーザーが不意に結末を知ってしまうことを防ぐシステムの基礎的検討を行う．
\end{abstract}

\section*{演習1：動作確認}
GitHubからこの「\texttt{LaTeX\_practice.tex}」をダウンロードし，
自分のPCでコンパイルしてみましょう．
正しくコンパイルできれば，PDFが生成されます．

\section*{演習2：テーマの変更}
タイトルと著者名を自分のものに変更し，\texttt{date}には\verb|\today|を入れて，
コンパイルしてみましょう．

\section*{演習3：Abstractの挿入}
Abstractを挿入することで，論文の概要を簡潔にまとめることができます．

\begin{verbatim}
\begin{abstract}
ここにAbstractを記述します．
\end{abstract}
\end{verbatim}

実際に，このコードを\verb|\maketitle|の後へ挿入してみましょう．

\section*{演習4：数式の挿入}
LaTeXでは，数式を美しく記述することができます．
数式の記述には，インライン数式とディスプレイ数式の2種類があります．

\subsection*{インライン数式}
円の面積は $S=\pi r^2$ で表される．

オイラーの公式は $e^{i\pi}+1=0$ である．

関数 $f(x)=x^2+2x+1$ を考える．

\subsection*{ディスプレイ数式}
\begin{equation}
\int_{0}^{1} x^2\,\dif x=\frac{1}{3}
\label{eq:integral}
\end{equation}

\begin{equation}
A=
\begin{pmatrix}
1 & 2\\
3 & 4
\end{pmatrix}
\label{eq:matrix}
\end{equation}

\begin{equation}
\lim_{n\to\infty}
\left(1+\frac{1}{n}\right)^n=e
\label{eq:euler}
\end{equation}

ディスプレイ数式は，\verb|\[ ... \]|で囲むと番号なしで表示できます．
番号を付けたい場合は，\verb|equation|環境を使用します．
実際に，インライン数式とディスプレイ数式をそれぞれ使って，
自分で数式を挿入してみましょう．

出力層でよく用いられるソフトマックス関数は，次のように表される．
\begin{equation}
y_k = \frac{\exp(a_k)}{\sum_{i=1}^{n} \exp(a_i)}
\label{eq:softmax}
\end{equation}

\section*{演習5：図の挿入}
LaTeXでは，図を挿入することもできます．

\begin{figure}[H]
    \centering
    \practiceimage[3cm]{figures/LaTeX_practice_figure_1.jpg}
    \caption{LaTeXで挿入した図の例}
    \label{fig:practice}
\end{figure}

画像の挿入は，画像ファイルの場所や拡張子，パッケージ設定などが原因で
エラーになることがあります．
エラーが出ても慌てず，エラーメッセージを読みながら一つずつ原因を確認しましょう．
では，自分で画像を挿入してみましょう．

\begin{figure}[H]
    \centering
    % もし figures フォルダの中ではなく、一番上の階層にアップロードした場合は
    % {nyanso.png} にしてください。
    \practiceimage[5cm]{figures/nyanso.png.PDF}
    \caption{アップロードした画像}
    \label{fig:nyanso}
\end{figure}

\section*{演習6：表の挿入}
LaTeXでは，表を挿入することもできます．

\begin{table}[H]
    \centering
    \begin{tabular}{|c|c|c|}
        \hline
        A & B & C \\
        \hline
        1 & 2 & 3 \\
        4 & 5 & 6 \\
        \hline
    \end{tabular}
    \caption{LaTeXで挿入した表の例}
    \label{tab:practice}
\end{table}

では，自分で5行7列の表を挿入してみましょう．

\begin{table}[H]
    \centering
    \caption{各モデルのハイパーパラメータと評価指標の比較}
    \label{tab:model_comparison}
    \resizebox{\columnwidth}{!}{% 幅がはみ出ないように調整
    \begin{tabular}{|l|c|c|c|c|c|c|}
        \hline
        モデル & Epochs & Batch Size & 精度 & 再現率 & F1スコア & 実行時間 \\
        \hline
        Model A & 10 & 32 & 0.85 & 0.82 & 0.83 & 12.5s \\
        Model B & 20 & 64 & 0.89 & 0.87 & 0.88 & 24.1s \\
        Model C & 30 & 128 & 0.92 & 0.91 & 0.91 & 35.8s \\
        Model D & 50 & 256 & 0.94 & 0.93 & 0.93 & 58.2s \\
        \hline
    \end{tabular}
    }
\end{table}

\section*{演習7：箇条書き}
\begin{itemize}
    \item LaTeX
    \item Python
    \item GitHub
\end{itemize}

では，自分で箇条書きを作成してみましょう．

自然言語処理において重要な要素は以下の通りである．
\begin{itemize}
    \item 形態素解析
    \item 単語のベクトル化 (Word Embedding)
    \item 文脈の理解と分類モデルの構築
\end{itemize}

\section*{演習8：ソースコードの挿入}
LaTeXでは，ソースコードもきれいに掲載できます．

\begin{lstlisting}
#include <stdio.h>

int main(void)
{
    for (int i = 0; i < 10; i++) {
        printf("%d\n", i);
    }

    return 0;
}
\end{lstlisting}

では，自分でソースコードを挿入してみましょう．

\begin{lstlisting}[language=Python]
import numpy as np

def sigmoid(x):
    """シグモイド関数"""
    return 1 / (1 + np.exp(-x))

# サンプルデータでの動作確認
x = np.array([-1.0, 0.0, 1.0])
print(sigmoid(x))
\end{lstlisting}

\section*{演習9：URLの挿入}
URLを挿入することもできます．
OpenAIのホームページはこちらです．

\url{https://openai.com}

では，自分で任意のURLを挿入してみましょう．

機械学習モデルの構築には，Hugging Faceのライブラリが非常に有用である．\\
\url{https://huggingface.co/}

\section*{演習10：参考文献}
卒業論文や学術論文では，他者の研究や論文を引用した場合，
必ず参考文献を明記する必要があります．

LaTeXでは，一般的にBibTeXを用いて参考文献を管理します．
参考文献を追加する手順は，大きく分けて次の3つです．

\subsection*{手順1：\texttt{.bib}ファイルを作成する}
まず，\texttt{references.bib}という名前のファイルを作成します．
このファイルには，引用する論文や書籍の情報を記述します．
例えば，次のように記述します．

\begin{lstlisting}[language={}]
@inproceedings{he2016resnet,
  author    = {Kaiming He and others},
  title     = {Deep Residual Learning for Image Recognition},
  booktitle = {CVPR},
  year      = {2016}
}
\end{lstlisting}

Google ScholarやIEEE Xploreでは，「BibTeX」を選択するだけで
この形式を取得できるため，自分で一から入力する必要はほとんどありません．

\subsection*{手順2：本文中で引用する}
引用したい箇所に，次のように記述します．

\begin{verbatim}
\cite{he2016resnet}
\end{verbatim}

例えば，

\begin{verbatim}
深層学習ではResNetが広く利用されている
\cite{he2016resnet}．
\end{verbatim}

と書くと，「深層学習ではResNetが広く利用されている[1]．」のように，
引用番号が自動で付与されます．

\subsection*{手順3：参考文献一覧を表示する}
文書の末尾へ，次の2行を追加します．

\begin{verbatim}
\bibliographystyle{ieeetr}
\bibliography{references}
\end{verbatim}

ここで，

\begin{itemize}
    \item \texttt{ieeetr}：参考文献の表示形式
    \item \texttt{references}：作成した\texttt{references.bib}の指定
\end{itemize}

を表しています．
実際に手順1から手順3までを行うと，本文では次のように表示されます．

「深層学習ではResNetが広く利用されている\cite{he2016resnet}．」

卒業論文では，参考文献の管理にBibTeXを用いることが一般的です．
今回は演習を省略しますが，研究を進める際には必ず利用する機能です．

\subsection*{発展：2段組で大きな図を挿入する}
卒業論文では，本文を2段組にし，図やグラフのみページ全体の幅で
表示することがあります．

その場合は，\verb|figure|の代わりに\verb|figure*|環境を使用します．
図がページ上部へ自動配置されることを確認するため，後ろにダミーテキストを入れています．

\begin{lstlisting}[language={[LaTeX]TeX}]
\begin{figure*}[t]
    \centering
    \includegraphics[width=0.8\textwidth]{graph.pdf}
    \caption{実験結果}
    \label{fig:result}
\end{figure*}
\end{lstlisting}

\begin{figure*}[t]
    \centering
    \begin{subfigure}{0.47\textwidth}
        \centering
        \practiceimage[\linewidth]{figures/LaTeX_practice_figure_1.jpg}
        \caption{画像1}
        \label{fig:image1}
    \end{subfigure}
    \hfill
    \begin{subfigure}{0.47\textwidth}
        \centering
        \practiceimage[\linewidth]{figures/LaTeX_practice_figure_2.jpg}
        \caption{画像2}
        \label{fig:image2}
    \end{subfigure}
    \caption{画像を横並びにした例}
    \label{fig:images}
\end{figure*}

\lipsum[2-7]

% 参考文献一覧は，文書の最後に一度だけ出力する．
\clearpage
\bibliographystyle{ieeetr}
\bibliography{references}

\end{document}